\begin{definition}
    Говорят, что на вероятностном пространстве $(\Omega, \B, P)$ задана \textit{фильтрация}, если задано семейство под-$\sigma$-алгебр $\F_t$ в $\B$, индексируемых параметром $t$ из некоего множества $T$ на прямой (или иного упорядоченного множества), причем $\F_s \subset \F_t$ при $s \leq t$.
\end{definition}

\begin{definition}
    Случайный процесс $\{\xi_t\}_{t \in T}$ называется \textit{согласованным 
    с фильтрацией} $\{\F_t\}_{t \in T}$, если при каждом $t$ случайная величина 
    $\xi_ t$ измерима относительно $\F_t$.
\end{definition}

\begin{note}
    Всякий случайный процесс $\{\xi_t\}_{t \in T}$ согласован с фильтрацией, 
    им порожденной, когда в качестве $\F_t$ берется $\sigma$-алгебра, порожденная 
    всеми случайными величинами $\xi_s$ с $s \leq t$. Для обсуждения марковских процессов такую $\sigma$-алгебру удобно обозначать символом $\F_{\leq t}$.
\end{note}

\begin{definition}
    Случайный процесс $\{\xi_t\}_{t \in T}$ называется \textit{мартингалом}
    относительно фильтрации $\{\F_t\}_{t \in T}$, если он с ней согласован, случайные величины $\xi_t$ интегрируемы и при $s \leq t$ п.в. верно равенство

    \[
    \E(\xi_t|\F_s) = \xi_s
    \]

    Если вместо равенства выполняется неравенство $\E(\xi_t|\F_s) \geq \xi_s$, то 
    процесс называется \textit{субмартингалом} относительно $\{\F_t\}_{t \in T}$.
    Если неравенство выполнено в другую сторону ($\E(\xi_t|\F_s) \leq \xi_s$), то
    процесс называется \textit{супермартингалом} относительно $\{\F_t\}_{t \in T}$.
\end{definition}

\begin{note}
    \(\E \xi_t = \E \xi_s\), так как \(\E \xi_t = \E(\E(\xi_t|\F_s)) = \E \xi_s\).
    Если $T \subset [0, +\infty)$, $0 \in T$ и $\xi_0 = 0$ или хотя бы $\E \xi_0 = 0$, то $\E \xi_t = 0$ для всех $t$.
\end{note}

\begin{example}[мартингала]
    Пусть $S_n = \xi_1 + \ldots + \xi_n$,
    где $\xi_i$ независимы и имеют нулевые математические ожидания, $T = \N$, 
    $\F_n$ порождена $\xi_1, \ldots, \xi_n$. Тогда $S_n$ является мартингалом. 
    % В самом деле, при $m > n$ имеет место равенство $\E(\xi_m|\F_n) = \E \xi_m = 0$ в силу задачи 11 выше.
\end{example}

\begin{proof}
    Пусть $k < n$. Надо показать, что $\E(S_n | \F_k) = S_k$. Достаточно доказать для $k = n - 1$.
    \[
    \E(S_n | \F_{n - 1}) = 
    \underbrace{\E (\xi_1 + \ldots + \xi_{n - 1} | \F_{n - 1})}_{\text{все эти функции измеримы относительно } \F_{n - 1}}
    + \E(\xi_n | \F_{n - 1}) = S_{n - 1} + \E(\xi_n | \F_{n - 1})
    \]

    Надо показать, что $\E(\xi_n | \F_{n - 1}) = 0$. Это эквивалентно

    \[
    \int_A \xi_n dP = 0 \ \forall A \in \sigma(\xi_1, \ldots, \xi_{n - 1})
    \]

    $I_A$ устроен следующим образом: $I_A = \phi(\xi_1, \ldots, \xi_{n - 1})$, 
    где $\phi$~--- борелевская на $\R^{n - 1}$.
    Тогда интеграл можно переписать в виде:

    \[
    \int_A \xi_n dP = \int_{\Omega} \phi(\xi_1, \ldots, \xi_{n - 1}) \xi_n dP
    \stackrel{\text{независимость}}{=}
    \int_{\Omega} \phi(\xi_1, \ldots, \xi_{n - 1}) dP \cdot
    \underbrace{\int_{\Omega} \xi_n dP}_{ = \E \xi_n = 0} = 0
    \]
\end{proof}

\begin{theorem}[Дуба о сходимости мартингалов]
    Пусть $\{\xi_n\}_{n \in \N}$~--- мартингал относительно $\{\F_n\}_{n \in \N}$. 
    Этот мартингал сходится в $L_1(P)$ в точности тогда, когда есть такая функция $\xi \in L_1(P)$, что $\xi_n$ = $\E^{\F_n} \xi$. 
    При этом есть и сходимость почти всюду.

    Равносильным условием является равномерная интегрируемость $\xi_n$, т.е. условие

    \[
    \lim_{R \to +\infty} \sup_n \int_{\{|\xi_n| \geq R\}} |\xi_n| dP = 0
    \]
\end{theorem}

\begin{theorem}[Критерий мартингальности для процессов с независимыми приращениями]
    Пусть $\{\xi_t\}$~--- процесс с независимыми приращениями. Он является мартингалом тогда и только тогда, когда $\E \xi_t = const$ ($\E \xi_t = \E \xi_s)$.
\end{theorem}

\begin{proof}
    \[
    \E(\xi_t | \F_s) = \E(\xi_t - \xi_s + \xi_s | \F_s) = 
    \underbrace{\E(\xi_t - \xi_s | \F_s)}_{\xi_t - \xi_s \indep \F_s} +
    \underbrace{\E(\xi_s | \F_s)}_{\xi_s \text{ измерима относительно } \F_s} =
    \E(\xi_t - \xi_s) + \xi_s
    \]
    
    $\{\xi_t\}$~--- мартингал $\lra \E(\xi_t | \F_s) = X_s \lra \E(\xi_t - \xi_s) = 0 \lra \E \xi_t = \E \xi_s$ 
\end{proof}

\begin{corollary}
    \noindent
    \begin{itemize}
        \item Винеровский процесс $(W_t, t \geq 0)$~--- мартингал
        \item Пуассоновский процесс $(N_t, t \geq 0)$~--- субмартингал
    \end{itemize}    
\end{corollary}
